\documentclass[12pt]{article}
\usepackage[utf8]{inputenc}
\usepackage[T1]{fontenc}
\usepackage[margin=1in]{geometry}
\usepackage{setspace}
\usepackage{xcolor}
\usepackage[colorlinks=true, urlcolor=blue!60!black, linkcolor=blue!60!black]{hyperref}
\setlength{\parindent}{0em}
\setlength{\parskip}{0.6em}
\renewcommand{\baselinestretch}{1.0}
\begin{document}
\begin{center}
{\LARGE \textit{Kelo v. City of New London}, 545 U.S. 469 (2005)\par}
\vspace{0.5em}
{\large Case Brief --- Two-Page Edition\par}
\end{center}

\section{Memory Jogger}
Economic development is a ``public purpose'': New London could condemn unblighted homes under a comprehensive plan, though States may set stricter limits (\href{https://supreme.justia.com/cases/federal/us/545/469/}{Justia}).

\section{Detailed Case Facts}
Economically distressed New London adopted a 90-acre Fort Trumbull redevelopment plan, next to a new Pfizer facility, with a hotel, offices, housing, and a marina. After most owners sold, the city condemned 15 well-kept, unblighted properties, including Susette Kelo's renovated home and the house where Wilhelmina Dery had lived since 1918.

\section{Procedural History}
The trial court upheld part of the takings; the Connecticut Supreme Court, 4--3, upheld them all under \textit{Berman} and \textit{Midkiff}.

\section{Judicial Votes}
Affirmed, 5--4. Stevens wrote for Kennedy, Souter, Ginsburg, and Breyer; Kennedy concurred; O'Connor, with Rehnquist, Scalia, and Thomas, dissented; Thomas also dissented separately.

\section{Holding}
Takings for transfer to private parties under a carefully formulated economic-development plan serve a public purpose and satisfy the Public Use Clause under deferential review, even without blight or public access.

\section{Analysis of Opinions}
\textbf{Majority:} The Court long ago rejected a literal ``use by the public'' test; plans are judged as a whole; there was no ``mere pretext,'' and courts will not demand ``reasonable certainty'' of benefits. \textbf{Kennedy:} takings favoring a particular private party with only pretextual public benefits remain invalid. \textbf{O'Connor:} the ruling washes out the line between public and private use. \textbf{Thomas:} overrule \textit{Berman} and \textit{Midkiff}; deference harms the powerless.

\section{Examples: Future Applications}
\textbf{Same-Side:} a county's studied, 200-acre industrial-park plan with tenants chosen later is valid. \textbf{Opposite-Side:} condemning a single nursery for a corporation that failed to buy it, with no plan, is likely pretextual. \textbf{Fence-Sitter:} a downtown plan built around a named team's arena.

\section{Critique}
Fair-market compensation omits subjective value and relocation costs, enriching private transferees, and Kennedy's pretext test is hard to apply when private profit and public tax gains are intertwined.

\section{Key Quotations}
\begin{quote}
``Promoting economic development is a traditional and long accepted function of government. There is, moreover, no principled way of distinguishing economic development from the other public purposes that we have recognized. \dots\ We emphasize that nothing in our opinion precludes any State from placing further restrictions on its exercise of the takings power.''
\end{quote}
\noindent --- Justice Stevens, Opinion of the Court, 545 U.S.\ at 484, 489.

\end{document}
